% ch4.tex -- Module 6, Chapter 4: Transformers and Power Transmission
% Dot points: M6.IQ3.03, M6.IQ3.04, M6.IQ3.05.
% (Faraday's iron ring / two solenoids, M6.IQ3.02, is a key experiment in Chapter 3.)
% All figures are TikZ / pgfplots, drawn inline.
\chapter{Transformers and Power Transmission}\label{ch4:chap}
\covers{M6.IQ3.03, M6.IQ3.04, M6.IQ3.05}

% chapter-local colours (prefix chFour, as the style guide asks)
\colorlet{chFourcopper}{hscseriesB!80!black}  % coil wire
\colorlet{chFouriron}{hscmuted!28}            % iron core
\colorlet{chFourcurrent}{hscseriesB}          % induced (eddy) currents

\inquiry{How are electric and magnetic fields related?}

\begin{goals}
  \item explain how a transformer uses a changing flux to change an AC voltage
  \item use $\dfrac{V_\mathrm{p}}{V_\mathrm{s}} = \dfrac{N_\mathrm{p}}{N_\mathrm{s}}$ and
        $V_\mathrm{p}I_\mathrm{p} = V_\mathrm{s}I_\mathrm{s}$ for ideal transformers
  \item evaluate why real transformers waste energy, and how to reduce the losses
  \item explain why electricity is sent at high voltage, using $P = VI$ and
        $P_\mathrm{loss} = I^2R$
\end{goals}

\begin{syllabus}
  \item M6.IQ3.03: analyse with numbers how ideal transformers work, using
        $\dfrac{V_\mathrm{p}}{V_\mathrm{s}} = \dfrac{N_\mathrm{p}}{N_\mathrm{s}}$ and
        $V_\mathrm{p}I_\mathrm{p} = V_\mathrm{s}I_\mathrm{s}$
  \item M6.IQ3.04: evaluate the limits of the ideal transformer model, and how
        transformers are made more efficient: incomplete flux linkage, resistive heating
        and eddy currents
  \item M6.IQ3.05: analyse uses of step-up and step-down transformers, including sending
        energy through high-voltage transmission lines
\end{syllabus}
Every phone charger and every power station uses a transformer. Transformers let us send
electricity hundreds of kilometres with little loss. Then they make it safe for your home.


%% ==================================================================
\section{The ideal transformer}\label{ch4:sec:ideal}
\covers{M6.IQ3.03}

\begin{bigidea}
A transformer changes the size of an AC voltage. A changing current in one coil makes a
changing magnetic field in an iron core. This changing field makes (induces) a voltage in
a second coil. Energy is conserved, so when the voltage goes up, the current goes down.
\end{bigidea}

\begin{keywords}
  \kw{Direct current (DC)}{current that flows in one direction only.}
  \kw{Alternating current (AC)}{current that keeps reversing its direction.}
  \kw{Transformer}{a device that changes the size of an AC voltage.}
  \kw{Primary / secondary coil}{the input coil (joined to the supply) / the output coil
    (joined to the load, the device that uses the energy).}
  \kw{Turns ratio}{$N_\mathrm{p} : N_\mathrm{s}$, primary turns compared with secondary
    turns.}
  \kw{Mutual induction}{an emf made in one coil by a changing current in a nearby coil.}
  \kw{Ideal transformer}{a perfect transformer that wastes no energy.}
  \kw{Step-up / step-down}{makes the voltage bigger / smaller.}
\end{keywords}

\subsection{DC and AC}

\begin{sidebyside}[0.46]
  \centering
  \begin{tikzpicture}
    \begin{axis}[hsc, width=66mm, height=42mm, xmin=0, xmax=40, ymin=-3.6, ymax=4.6,
        xtick={0,10,20,30,40}, ytick={-2,0,2},
        xlabel={Time $t$ (\si{ms})}, ylabel={Current $I$ (\si{A})}]
      \draw[hscink!60, line width=0.5pt] (axis cs:0,0) -- (axis cs:40,0);
      \addplot[hscseriesA, line width=1.2pt, domain=0:40, samples=2] {1.5};
      \addplot[hscseriesB, line width=1.2pt, domain=0:40, samples=161, smooth] {2.5*sin(18*x)};
      \node[font=\small\bfseries, text=hscseriesA, anchor=south] at (axis cs:36.5,1.5) {DC};
      \node[font=\small\bfseries, text=hscseriesB, anchor=north] at (axis cs:15,-2.6) {AC};
      \draw[guide] (axis cs:5,2.5) -- (axis cs:5,3.5) (axis cs:25,2.5) -- (axis cs:25,3.5);
      \draw[<->, hscink, line width=0.6pt] (axis cs:5,3.3) -- (axis cs:25,3.3)
        node[midway, above, inner sep=1.5pt, font=\footnotesize] {one period $T = \qty{20}{ms}$};
    \end{axis}
  \end{tikzpicture}
  \captionof{figure}{Current against time for DC (blue) and for \qty{50}{Hz} AC
    (orange).}
  \label{ch4:fig:dcac}
\nextside
  A battery gives \term{direct current} (DC): current that flows in one direction only.
  The power points (wall sockets) at home give \term{alternating current} (AC): current
  that keeps reversing its direction (\cref{ch4:fig:dcac}).\par\smallskip
  Australian \term{mains} electricity (the supply from the wall socket) is \qty{230}{V}
  AC. Older books say \qty{240}{V}. If a question gives \qty{240}{V}, use \qty{240}{V}.
\end{sidebyside}

One \term{cycle} is one full swing forwards and back. The \term{frequency} ($f$) is the
number of cycles each second; for mains it is \qty{50}{Hz}. The \term{period} ($T$) is the
time for one cycle. The sheet gives $f = \dfrac{1}{T}$, so
$T = \dfrac{1}{f} = \dfrac{1}{50} = \qty{0.020}{s} = \qty{20}{ms}$
(\qty{1}{ms} $=$ \qty{0.001}{s}).

A current makes a magnetic field. AC keeps changing, so its magnetic field keeps changing
too. A changing field is what Faraday's law needs to induce a voltage (see the refresher
below).

\begin{refresher}{Fields from currents, Faraday's law and Lenz's law (\cref{ch3:chap})}
\begin{itemize}
  \item A current in a coil makes a magnetic field through the coil
        (\cref{ch2:sec:fields}). A bigger current makes a stronger field. If the current
        reverses, the field reverses too.
  \item \term{Magnetic flux} ($\Phi = BA\cos\theta$, unit \si{Wb}) is how much magnetic
        field passes through an area.
  \item Faraday's law: a \emph{changing} flux through a coil induces an \term{emf} (a
        voltage made by induction), $\varepsilon = -N\dfrac{\Delta\Phi}{\Delta t}$. No
        change in flux, no emf. The minus sign is Lenz's law: the induced current makes a
        magnetic field that opposes the change in flux.
\end{itemize}
\end{refresher}

\subsection{How a transformer works}

Your phone needs about \qty{5}{V}, but the wall socket gives \qty{230}{V}. A transformer
inside the charger makes the voltage smaller.

A \term{transformer} is a device that changes the size of an AC voltage.
\Cref{ch4:fig:transformer} shows a simple one. The \term{primary coil}
($N_\mathrm{p}$ turns) is joined to the AC supply. The \term{secondary coil}
($N_\mathrm{s}$ turns) is joined to the \term{load}: the device that uses the energy.
One \term{turn} is one loop of wire around the core. Both coils are wrapped around one
closed \term{core}: a ring of iron that carries the magnetic field from one coil to the
other. The core is made of \term{soft iron}, which is iron that is easy to turn into a
magnet and back again. No wire joins the two coils.

\begin{figure}
  \centering
  \begin{tikzpicture}[x=1cm, y=1cm]
    % iron core (closed loop)
    \fill[chFouriron, even odd rule] (0,0) rectangle (5.6,3.6) (1.1,0.9) rectangle (4.5,2.7);
    \draw[hscink!70, line width=0.6pt] (0,0) rectangle (5.6,3.6) (1.1,0.9) rectangle (4.5,2.7);
    % flux path (clockwise at this instant)
    \draw[hscfield, line width=0.8pt, dashed] (0.55,0.45) rectangle (5.05,3.15);
    \draw[field] (1.4,3.15) -- (4.2,3.15);
    \draw[field] (4.2,0.45) -- (1.4,0.45);
    \node[diagramlabel, text=hscfield, align=center] at (2.8,1.8)
      {changing\\ magnetic flux $\Phi$\\ in the core};
    % primary coil: 4 turns on the left limb
    \foreach \y in {1.2,1.6,2.0,2.4}
      \draw[chFourcopper, line width=1.6pt] (-0.14,\y) -- (1.24,\y+0.18);
    % secondary coil: 8 turns on the right limb
    \foreach \y in {0.95,1.17,...,2.5}
      \draw[chFourcopper, line width=1.6pt] (5.74,\y) -- (4.36,\y+0.11);
    % primary circuit and AC supply (leads join the two ends of the coil)
    \draw[hscink, line width=0.7pt] (-0.14,2.4) -- (-2.2,2.4) -- (-2.2,2.15);
    \draw[hscink, line width=0.7pt] (-0.14,1.2) -- (-2.2,1.2) -- (-2.2,1.45);
    \draw[hscink, line width=0.7pt, fill=white] (-2.2,1.8) circle (0.35);
    \node at (-2.2,1.8) {$\sim$};
    \node[diagramlabel, align=center, below] at (-2.2,1.1) {AC\\ supply};
    \node[diagramlabel] at (-1.2,1.8) {$V_\mathrm{p}$};
    \node[diagramlabel, above] at (-1.2,2.4) {$I_\mathrm{p}$};
    % secondary circuit and lamp (leads join the two ends of the coil)
    \draw[hscink, line width=0.7pt] (5.74,2.49) -- (7.6,2.49) -- (7.6,2.1);
    \draw[hscink, line width=0.7pt] (5.74,0.95) -- (7.6,0.95) -- (7.6,1.5);
    \draw[hscink, line width=0.7pt, fill=white] (7.6,1.8) circle (0.3);
    \draw[hscink, line width=0.7pt] (7.6,1.8) ++(45:0.3) -- ++(225:0.6)
      (7.6,1.8) ++(135:0.3) -- ++(315:0.6);
    \node[diagramlabel, align=center, below] at (7.6,0.85) {lamp\\ (load)};
    \node[diagramlabel] at (6.7,1.75) {$V_\mathrm{s}$};
    \node[diagramlabel, above] at (6.7,2.49) {$I_\mathrm{s}$};
    % labels
    \node[diagramlabel] (pl) at (0.3,4.0) {primary coil ($N_\mathrm{p}$ turns)};
    \draw[hscmuted, line width=0.5pt] (0.3,3.8) -- (0.3,2.52);
    \node[diagramlabel] (sl) at (5.3,4.0) {secondary coil ($N_\mathrm{s}$ turns)};
    \draw[hscmuted, line width=0.5pt] (5.3,3.8) -- (5.3,2.58);
    \node[diagramlabel, below] at (2.8,-0.05) {soft-iron core (a closed loop)};
  \end{tikzpicture}
  \caption{A step-up transformer ($N_\mathrm{s} > N_\mathrm{p}$). The flux arrows show one
    instant; with AC they reverse every half-cycle. No wire joins the two coils.}
  \label{ch4:fig:transformer}
\end{figure}

Energy goes from the primary circuit to the secondary circuit through the changing
magnetic field. Here is how, one step at a time.
\begin{enumerate}
  \item AC in the primary makes a magnetic field that keeps changing size and direction.
  \item The iron core carries this changing flux around to the secondary coil.
  \item The flux through each secondary turn keeps changing, so an emf is induced in
        every turn (Faraday's law).
  \item The turns are in series (joined one after another along the same wire), so
        their emfs add up: more turns, bigger voltage.
  \item The flux reverses with the supply, so the output is AC of the \emph{same
        frequency}.
\end{enumerate}
An emf made in one coil by a changing current in a nearby coil is called
\term{mutual induction}.

\textbf{Why a transformer needs AC.} Steady DC also makes a magnetic field, but the flux
does not change. So $\dfrac{\Delta\Phi}{\Delta t} = 0$ and no emf is induced. The
secondary gives only a short pulse when the DC is switched on or off (Faraday's iron
ring, \cref{ch3:sec:experiments}).

\textbf{Step-up and step-down.} A \term{step-up transformer} has more turns on the
secondary ($N_\mathrm{s} > N_\mathrm{p}$), so it makes the voltage bigger. A
\term{step-down transformer} has fewer ($N_\mathrm{s} < N_\mathrm{p}$), so it makes the
voltage smaller. The \term{turns ratio} $N_\mathrm{p} : N_\mathrm{s}$ compares the turns.

\begin{toolkit}{Ratios and rearranging a fraction equation}
``$N_\mathrm{p} : N_\mathrm{s} = 1 : 20$'' means the secondary has 20 times as many turns
as the primary. To find $V_\mathrm{s}$ in $\dfrac{12}{V_\mathrm{s}} = \dfrac{100}{400}$:
cross-multiply (multiply each top by the bottom of the other fraction):
$12 \times 400 = 100\,V_\mathrm{s}$; then divide both sides by 100,
$V_\mathrm{s} = \dfrac{12 \times 400}{100} = \qty{48}{V}$.
\end{toolkit}

\begin{formulacard}{The transformer (turns) equation}
  \cardformula{\frac{V_\mathrm{p}}{V_\mathrm{s}} = \frac{N_\mathrm{p}}{N_\mathrm{s}}}
  \begin{symbols}
    \sym{$V_\mathrm{p}$}{voltage across the primary coil (input)}{\si{V}}
    \sym{$V_\mathrm{s}$}{voltage across the secondary coil (output)}{\si{V}}
    \sym{$N_\mathrm{p}$, $N_\mathrm{s}$}{number of turns on the primary, on the
      secondary}{no unit}
  \end{symbols}
  \inwords{the voltage ratio equals the turns ratio. More secondary turns give a bigger
    output voltage.}
  \whentouse{a transformer on AC (exact if ideal, a good estimate if real).}
  \whennottouse{steady DC (the output is zero), or for currents.}
  \rearranged{$V_\mathrm{s} = V_\mathrm{p}\dfrac{N_\mathrm{s}}{N_\mathrm{p}}$ \qquad
    $N_\mathrm{s} = N_\mathrm{p}\dfrac{V_\mathrm{s}}{V_\mathrm{p}}$ \qquad
    $N_\mathrm{p} = N_\mathrm{s}\dfrac{V_\mathrm{p}}{V_\mathrm{s}}$}
  \tensecond{$N_\mathrm{p} = 500$, $N_\mathrm{s} = 50$, $V_\mathrm{p} = \qty{230}{V}$, so
    $V_\mathrm{s} = 230 \times \dfrac{50}{500} = \qty{23}{V}$.}
  \derive{the core carries the same flux through every turn of both coils, so each turn
    has the same $\dfrac{\Delta\Phi}{\Delta t}$. Use the size of Faraday's emf for each coil.
    \begin{align*}
      &V_\mathrm{s} = N_\mathrm{s}\frac{\Delta\Phi}{\Delta t}
        && \text{Faraday's law (size only) for the secondary}\\
      &V_\mathrm{p} = N_\mathrm{p}\frac{\Delta\Phi}{\Delta t}
        && \text{the same for the primary (in an ideal coil,}\\
      & && \text{\quad the emf induced in it equals the supply voltage)}\\
      &\frac{V_\mathrm{p}}{V_\mathrm{s}}
        = \frac{N_\mathrm{p}\,\Delta\Phi/\Delta t}{N_\mathrm{s}\,\Delta\Phi/\Delta t}
        = \frac{N_\mathrm{p}}{N_\mathrm{s}}
        && \text{divide line 2 by line 1; $\tfrac{\Delta\Phi}{\Delta t}$ cancels}
    \end{align*}
    This comes from Faraday's law, not from conservation of energy.}
  \trap{put the primary values on the top of both fractions ($V_\mathrm{p}$ over
    $V_\mathrm{s}$, and $N_\mathrm{p}$ over $N_\mathrm{s}$). The primary is the coil joined
    to the supply, wherever it is drawn.}
\end{formulacard}

\begin{sidebyside}[0.44]
  \textbf{What the output looks like.} The \term{peak} voltage is the largest voltage in
  each cycle. A transformer multiplies the peak by $\dfrac{N_\mathrm{s}}{N_\mathrm{p}}$.
  It does not change the period or the frequency (\cref{ch4:fig:vpvs}).

  With $N_\mathrm{p} = 600$ and $N_\mathrm{s} = 200$, a \qty{12}{V} input peak gives
  $12 \times \dfrac{200}{600} = \qty{4.0}{V}$.
\nextside
  \centering
  \begin{tikzpicture}
    \begin{axis}[hsc, width=76mm, height=46mm, xmin=0, xmax=40, ymin=-14, ymax=20,
        xtick={0,10,20,30,40}, ytick={-12,-4,4,12},
        xlabel={Time $t$ (\si{ms})}, ylabel={Voltage (\si{V})},
        legend style={at={(0.5,1.02)}, anchor=north, legend columns=2,
          /tikz/every even column/.append style={column sep=3mm}}]
      \draw[hscink!60, line width=0.5pt] (axis cs:0,0) -- (axis cs:40,0);
      \addplot[hscseriesA, line width=1.2pt, domain=0:40, samples=161, smooth] {12*sin(18*x)};
      \addplot[hscseriesB, line width=1.4pt, dashed, domain=0:40, samples=161, smooth] {4*sin(18*x)};
      \legend{input $V_\mathrm{p}$, output $V_\mathrm{s}$}
    \end{axis}
  \end{tikzpicture}
  \captionof{figure}{An ideal step-down transformer, \mbox{$600 : 200$}. The peak falls from
    \qty{12}{V} to \qty{4.0}{V}. The period (\qty{20}{ms}) stays the same.}
  \label{ch4:fig:vpvs}
\end{sidebyside}

The \qty{230}{V} of the mains is an \term{effective value}: the steady (DC) voltage that
would heat a heater by the same amount. The real voltage swings between about
$+\qty{325}{V}$ and $-\qty{325}{V}$ (its peak). The turns equation works with peak values
or with effective values, as long as both sides use the same kind.

\textbf{Careful: a graph of the primary \emph{current}.} Some questions give the primary
current (or the flux) against time, not the primary voltage. The flux in the core follows
the current. The secondary emf depends on how fast the flux changes
($\varepsilon = -N\dfrac{\Delta\Phi}{\Delta t}$). So the emf follows the \term{gradient}
of the current graph: how steep the graph is (\cref{ch4:fig:gradient}).
\begin{itemize}
  \item Where the current graph is flat, the emf is zero.
  \item Where it is steepest, the emf is biggest.
  \item A current that rises in a straight line gives a constant emf.
  \item A sine-shaped current gives a cosine-shaped emf, shifted by a quarter of a period.
\end{itemize}

\begin{figure}
  \centering
  \begin{tikzpicture}
    \begin{axis}[hsc, name=chFourItop, width=84mm, height=31mm, xmin=0, xmax=40,
        ymin=-1.35, ymax=1.35, xtick={0,10,20,30,40}, xticklabels={}, ytick=\empty,
        ylabel={$I_\mathrm{p}$}, ylabel style={rotate=-90}]
      \draw[hscink!60, line width=0.5pt] (axis cs:0,0) -- (axis cs:40,0);
      \draw[guide] (axis cs:5,-1.35) -- (axis cs:5,1.35) (axis cs:10,-1.35) -- (axis cs:10,1.35);
      \addplot[hscseriesA, line width=1.2pt, domain=0:40, samples=161, smooth] {sin(18*x)};
      \node[font=\footnotesize, text=hscseriesA, anchor=south] at (axis cs:5,1.0) {flat};
      \node[font=\footnotesize, text=hscseriesA, anchor=south west] at (axis cs:10.6,0.05) {steepest};
    \end{axis}
    \begin{axis}[hsc, at={(chFourItop.south)}, anchor=north, yshift=-3mm, width=84mm,
        height=31mm, xmin=0, xmax=40, ymin=-1.35, ymax=1.35, xtick={0,10,20,30,40},
        ytick=\empty, xlabel={Time $t$ (\si{ms})},
        ylabel={$\varepsilon_\mathrm{s}$}, ylabel style={rotate=-90}]
      \draw[hscink!60, line width=0.5pt] (axis cs:0,0) -- (axis cs:40,0);
      \draw[guide] (axis cs:5,-1.35) -- (axis cs:5,1.35) (axis cs:10,-1.35) -- (axis cs:10,1.35);
      \addplot[hscseriesB, line width=1.2pt, domain=0:40, samples=161, smooth] {-cos(18*x)};
      \node[font=\footnotesize, text=hscseriesB, anchor=north west] at (axis cs:5.4,-0.05) {zero};
      \node[font=\footnotesize, text=hscseriesB, anchor=south] at (axis cs:10,1.0) {biggest};
    \end{axis}
  \end{tikzpicture}
  \caption{Top: a sine-shaped primary current. Bottom: the secondary emf. It
    follows the gradient of the current graph, so it is shifted by a quarter of a period.}
  \label{ch4:fig:gradient}
\end{figure}

\subsection{Power and current}

A transformer cannot create energy. It only passes energy on. In an \term{ideal
transformer} no energy is wasted, so the power out equals the power in.

\begin{refresher}{Power and Ohm's law (Year 11)}
\term{Power} ($P$) is the energy transferred each second, in watts
($\qty{1}{W} = \qty{1}{J.s^{-1}}$). Ohm's law: $V = IR$, where $R$ is the resistance in
ohms (\si{\ohm}). Both are on the formula sheet.
\end{refresher}

\begin{formulacard}{Electrical power}
  \cardformula{P = VI}
  \begin{symbols}
    \sym{$P$}{power: energy transferred each second}{\si{W}}
    \sym{$V$}{voltage across that part of the circuit}{\si{V}}
    \sym{$I$}{current through that part of the circuit}{\si{A}}
  \end{symbols}
  \inwords{power equals voltage times current.}
  \whentouse{the power of one coil ($V_\mathrm{p}I_\mathrm{p}$ or
    $V_\mathrm{s}I_\mathrm{s}$), of a device, or the current a power station sends
    ($I = \dfrac{P}{V}$).}
  \whennottouse{to find the heat lost in a line from the voltage sent. Find $I$ first,
    then use $P_\mathrm{loss} = I^2R$ (\cref{ch4:sec:transmission}).}
  \rearranged{$I = \dfrac{P}{V}$ \qquad $V = \dfrac{P}{I}$}
  \tensecond{a kettle on \qty{230}{V} takes \qty{8.0}{A}, so
    $P = 230 \times 8.0 = \qty{1840}{W} \approx \qty{1.8}{kW}$.}
  \trap{$V$ and $I$ must belong to the same part of the circuit. Never multiply the
    primary voltage by the secondary current. Change \si{kV} to \si{V} and \si{mA} to
    \si{A} first.}
\end{formulacard}

\begin{formulacard}{Power in an ideal transformer}
  \cardformula{V_\mathrm{p} I_\mathrm{p} = V_\mathrm{s} I_\mathrm{s}}
  \begin{symbols}
    \sym{$V_\mathrm{p}$, $I_\mathrm{p}$}{voltage and current of the primary (input)}{\si{V}, \si{A}}
    \sym{$V_\mathrm{s}$, $I_\mathrm{s}$}{voltage and current of the secondary (output)}{\si{V}, \si{A}}
  \end{symbols}
  \inwords{power in equals power out. So when the voltage goes up, the current goes down
    by the same factor.}
  \whentouse{an ideal transformer, to find a current or a power.}
  \whennottouse{a transformer that loses energy (use the efficiency,
    \cref{ch4:sec:losses}).}
  \rearranged{$I_\mathrm{s} = \dfrac{V_\mathrm{p}I_\mathrm{p}}{V_\mathrm{s}}$ \qquad
    $I_\mathrm{p} = \dfrac{V_\mathrm{s}I_\mathrm{s}}{V_\mathrm{p}}$ \qquad and with the
    turns equation: $\dfrac{I_\mathrm{s}}{I_\mathrm{p}} = \dfrac{V_\mathrm{p}}{V_\mathrm{s}}
    = \dfrac{N_\mathrm{p}}{N_\mathrm{s}}$ (this form is not printed on the sheet)}
  \tensecond{$N_\mathrm{p} = 1000$, $N_\mathrm{s} = 50$, $I_\mathrm{p} = \qty{0.20}{A}$,
    so $I_\mathrm{s} = 0.20 \times \dfrac{1000}{50} = \qty{4.0}{A}$.}
  \derive{energy is conserved and nothing is lost, so power in = power out. With
    $P = VI$ for each coil: $V_\mathrm{p}I_\mathrm{p} = V_\mathrm{s}I_\mathrm{s}$. Divide
    both sides by $V_\mathrm{s}I_\mathrm{p}$: $\dfrac{V_\mathrm{p}}{V_\mathrm{s}} =
    \dfrac{I_\mathrm{s}}{I_\mathrm{p}}$. Then the turns equation gives the current form.}
  \trap{``step-up'' means the \emph{voltage} goes up; the current goes \emph{down}.
    $I_\mathrm{s}$ is on top in the current ratio.}
\end{formulacard}

\textbf{The load decides the current.} A \qty{60}{W} lamp on the secondary takes the
current it needs. The primary then draws just enough current to supply that power.

\begin{workedexample}{A garden-light transformer}\label{ch4:we:garden}
An ideal transformer changes \qty{230}{V} mains to \qty{12}{V} for garden lights. Its
primary has 1150 turns. The lights use \qty{60}{W} in total. Find (a)~the number of
secondary turns, (b)~the secondary current and (c)~the primary current.
\given{$V_\mathrm{p} = \qty{230}{V}$, \ $V_\mathrm{s} = \qty{12}{V}$, \
  $N_\mathrm{p} = 1150$, \ $P = \qty{60}{W}$ in and out (ideal)}
\find{(a) $N_\mathrm{s}$ \quad (b) $I_\mathrm{s}$ \quad (c) $I_\mathrm{p}$}
\formula{$N_\mathrm{s} = N_\mathrm{p}\dfrac{V_\mathrm{s}}{V_\mathrm{p}}$ \quad and \quad
  $P = VI$, so $I = \dfrac{P}{V}$ for each coil}
\substitute{(a) $N_\mathrm{s} = 1150 \times \dfrac{\qty{12}{V}}{\qty{230}{V}} = 60$
  \quad (the volts cancel)\par
  (b) $I_\mathrm{s} = \dfrac{\qty{60}{W}}{\qty{12}{V}} = \qty{5.0}{A}$ \qquad
  (c) $I_\mathrm{p} = \dfrac{\qty{60}{W}}{\qty{230}{V}} = \qty{0.261}{A}$}
\answer{(a) \final{60 turns} \quad (b) \final{$I_\mathrm{s} = \qty{5.0}{A}$} \quad
  (c) \final{$I_\mathrm{p} = \qty{0.26}{A}$}}
\checkstep{$\dfrac{I_\mathrm{s}}{I_\mathrm{p}} = \dfrac{5.0}{0.261} = 19.2$ and
  $\dfrac{N_\mathrm{p}}{N_\mathrm{s}} = \dfrac{1150}{60} = 19.2$ \cmark. Step-down, so
  the secondary has the bigger current \cmark.}
\end{workedexample}

\begin{workedexample}{The current trap}\label{ch4:we:trap}
An ideal transformer has 150 turns on its primary and 600 turns on its secondary. The
primary voltage is \qty{48}{V} and the secondary current is \qty{2.4}{A}.
(a)~Calculate the primary current. (b)~Calculate the output power.
\given{$N_\mathrm{p} = 150$, \ $N_\mathrm{s} = 600$, \ $V_\mathrm{p} = \qty{48}{V}$, \
  $I_\mathrm{s} = \qty{2.4}{A}$}
\find{(a) $I_\mathrm{p}$ \quad (b) $P_\mathrm{out} = V_\mathrm{s}I_\mathrm{s}$}
\formula{$\dfrac{I_\mathrm{s}}{I_\mathrm{p}} = \dfrac{N_\mathrm{p}}{N_\mathrm{s}}$
  (current ratio, from the formula card). Turn both fractions upside down:
  $\dfrac{I_\mathrm{p}}{I_\mathrm{s}} = \dfrac{N_\mathrm{s}}{N_\mathrm{p}}$. Multiply both
  sides by $I_\mathrm{s}$: $I_\mathrm{p} = I_\mathrm{s}\dfrac{N_\mathrm{s}}{N_\mathrm{p}}$.
  For (b): $V_\mathrm{s} = V_\mathrm{p}\dfrac{N_\mathrm{s}}{N_\mathrm{p}}$, then $P = VI$.}
\substitute{(a) $I_\mathrm{p} = \qty{2.4}{A} \times \dfrac{600}{150} = \qty{9.6}{A}$\par
  (b) $V_\mathrm{s} = \qty{48}{V} \times \dfrac{600}{150} = \qty{192}{V}$, so
  $P_\mathrm{out} = \qty{192}{V} \times \qty{2.4}{A} = \qty{460.8}{W}$}
\answer{(a) \final{$I_\mathrm{p} = \qty{9.6}{A}$} \quad
  (b) \final{$P_\mathrm{out} = \qty{4.6e2}{W}$} (2 significant figures)}
\checkstep{power in $= 48 \times 9.6 = \qty{460.8}{W}$ = power out \cmark. Step-up, so
  the primary has the \emph{bigger} current. (The upside-down ratio gives \qty{0.60}{A}.)}
\end{workedexample}

\begin{mistakes}
\begin{itemize}
  \item \xmark\ ``DC makes no magnetic field.'' \cmark\ DC makes a \emph{steady} field,
        so the flux does not change.
  \item \xmark\ $\dfrac{I_\mathrm{p}}{I_\mathrm{s}} = \dfrac{N_\mathrm{p}}{N_\mathrm{s}}$.
        \cmark\ The current ratio is upside down:
        $\dfrac{I_\mathrm{s}}{I_\mathrm{p}} = \dfrac{N_\mathrm{p}}{N_\mathrm{s}}$.
  \item \xmark\ ``A step-up transformer increases the power.''
        \cmark\ Power out is never more than power in.
  \item \xmark\ ``Current flows from one coil to the other.''
        \cmark\ No charge crosses; the changing field carries the energy.
\end{itemize}
\end{mistakes}

\begin{examtip}
``Explain how a transformer changes the voltage'' (3--4 marks). Write steps 1--4 of the
list in ``How a transformer works'' (\cref{ch4:sec:ideal}) as linked sentences. Markers
want:
\begin{enumerate}
  \item AC in the primary makes a changing magnetic flux;
  \item the iron core links this flux to the secondary;
  \item the changing flux through the secondary induces an emf (Faraday's law);
  \item $\dfrac{V_\mathrm{p}}{V_\mathrm{s}} = \dfrac{N_\mathrm{p}}{N_\mathrm{s}}$, so more
        secondary turns give a bigger voltage.
\end{enumerate}
In calculations, one mark is for the equation with numbers and one for the answer with a
unit. Graphs: from a graph of $V_\mathrm{p}$, keep the period and multiply the peak by
$N_\mathrm{s}/N_\mathrm{p}$. From a graph of the primary \emph{current}, draw its gradient
(\cref{ch4:fig:gradient}).
\end{examtip}

% Layout QA (whole book): the Quick check missed the bottom of the page by a few points
% and left 42 % of that page empty; one extra line on this page lets it fit.
\enlargethispage{\baselineskip}
\begin{quickcheck}
  \item A transformer with 400 primary turns and 2000 secondary turns is joined to
        \qty{24}{V} AC. Is it step-up or step-down? Find the output voltage.
  \item $N_\mathrm{p} = 250$ and $N_\mathrm{s} = 100$. The primary voltage has a peak of
        \qty{9.0}{V} and a frequency of \qty{50}{Hz}. State the peak and the frequency of
        the output.
  \item The primary current of a transformer rises steadily from 0 to \qty{2.0}{A} in
        \qty{5.0}{ms}. It stays at \qty{2.0}{A} for \qty{5.0}{ms}, then falls steadily to 0
        in \qty{5.0}{ms}. (a)~Sketch the secondary emf against time. (b)~Explain why the
        emf is zero from \qty{5}{ms} to \qty{10}{ms}.\label{ch4:qc:ramp}
  \item An ideal transformer has $N_\mathrm{p} : N_\mathrm{s} = 1 : 25$ and a primary
        current of \qty{50}{A}. Find the secondary current.
  \item A welding machine needs a very large current. Explain why it uses a step-down
        transformer.
\end{quickcheck}

%% ==================================================================
\section{Real transformers: energy losses and efficiency}\label{ch4:sec:losses}
\covers{M6.IQ3.04}

\begin{bigidea}
Real transformers waste a little energy, mostly as heat. So the power out is a bit less
than the power in. A closed iron core, thick wire and a laminated core keep the losses
small.
\end{bigidea}

\begin{keywords}
  \kw{Efficiency ($\eta$)}{the percentage of the input power that comes out as useful power.}
  \kw{Incomplete flux linkage}{some flux from the primary misses the secondary.}
  \kw{Resistive heating}{heat made when a current flows through a resistance.}
  \kw{Eddy currents}{loops of current induced inside a solid piece of metal by a changing
    magnetic flux.}
  \kw{Laminated core}{a core made of thin iron sheets with insulation between them.}
\end{keywords}

A phone charger feels warm after an hour of charging. That heat is energy that did not
reach the phone. In a real transformer the power out is less than the power in:
$V_\mathrm{s}I_\mathrm{s} < V_\mathrm{p}I_\mathrm{p}$. The energy is not destroyed. It
becomes heat (and a little sound). The syllabus names three main limitations of the
ideal model.

\subsection{Limitation 1: incomplete flux linkage}

\term{Incomplete flux linkage} means that not all the flux made by the primary passes
through the secondary. Some flux leaks into the air, so the secondary gets a smaller emf
than the turns equation predicts (\cref{ch4:fig:linkage}). Less energy reaches the
secondary.

\begin{figure}
  \centering
  \begin{tikzpicture}[x=0.85cm, y=0.85cm]
    % (a) two coils with no core
    \draw[field] (-1.55,0.12) -- (1.95,0.12) .. controls (2.9,0.12) and (2.9,1.3) ..
      (1.2,1.3) -- (-0.8,1.3) .. controls (-2.5,1.3) and (-2.5,0.12) .. (-1.55,0.12);
    \draw[field] (-1.55,-0.12) -- (1.95,-0.12) .. controls (2.9,-0.12) and (2.9,-1.3) ..
      (1.2,-1.3) -- (-0.8,-1.3) .. controls (-2.5,-1.3) and (-2.5,-0.12) .. (-1.55,-0.12);
    \draw[field] (-1.55,0.3) -- (-0.15,0.3) .. controls (0.55,0.3) and (0.45,0.85) ..
      (-0.8,0.85) .. controls (-2.0,0.85) and (-2.1,0.3) .. (-1.55,0.3);
    \draw[field] (-1.55,-0.3) -- (-0.15,-0.3) .. controls (0.55,-0.3) and (0.45,-0.85) ..
      (-0.8,-0.85) .. controls (-2.0,-0.85) and (-2.1,-0.3) .. (-1.55,-0.3);
    \foreach \x in {-1.4,-1.16,...,-0.3}
      \draw[chFourcopper, line width=1pt] (\x,0) ellipse (0.07 and 0.45);
    \foreach \x in {0.7,0.94,...,1.8}
      \draw[chFourcopper, line width=1pt] (\x,0) ellipse (0.07 and 0.45);
    \node[font=\small] at (-1.2,-1.7) {primary};
    \node[font=\small] at (1.6,-1.7) {secondary};
    \node[font=\small, text=hscfield] (lf) at (0.9,1.85) {leaked flux};
    \draw[hscmuted, line width=0.4pt] (lf.south west) -- (-0.1,0.72);
    \node[font=\small\bfseries, text=hscink] at (0.2,2.6) {(a) no core};
    % (b) closed iron core
    \begin{scope}[xshift=4.94cm, yshift=-1.33cm]
      \fill[chFouriron, even odd rule] (0,0) rectangle (3.4,2.5) (0.75,0.65) rectangle (2.65,1.85);
      \draw[hscink!70, line width=0.6pt] (0,0) rectangle (3.4,2.5) (0.75,0.65) rectangle (2.65,1.85);
      \draw[hscfield, line width=0.8pt, dashed] (0.375,0.325) rectangle (3.025,2.175);
      \draw[field] (0.9,2.175) -- (2.5,2.175);
      \draw[field] (2.5,0.325) -- (0.9,0.325);
      \foreach \y in {0.85,1.15,1.45}
        \draw[chFourcopper, line width=1.4pt] (-0.12,\y) -- (0.87,\y+0.14);
      \foreach \y in {0.75,0.95,...,1.6}
        \draw[chFourcopper, line width=1.4pt] (3.52,\y) -- (2.53,\y+0.1);
      \node[font=\small] at (0.375,-0.4) {primary};
      \node[font=\small] at (3.025,-0.4) {secondary};
      \node[font=\small\bfseries, text=hscink] at (1.7,3.65) {(b) closed iron core};
    \end{scope}
  \end{tikzpicture}
  \caption{(a) Without a core, much of the flux (violet) misses the secondary.
    (b) A closed iron core guides almost all of it through both coils.}
  \label{ch4:fig:linkage}
\end{figure}

How engineers reduce it:
\begin{itemize}
  \item wrap both coils around one \emph{closed} loop of soft iron, which guides the field
        lines from one coil to the other;
  \item wrap the secondary directly over the primary, on the same part of the core.
\end{itemize}
Wireless phone chargers have an air gap and no shared core, so much of their flux leaks.

\subsection{Limitation 2: resistive heating in the coils}

The coils are long copper wires, and every wire has some resistance. A current in a
resistance makes heat. This is \term{resistive heating}. The heat made each second is
$I^2R$, where $I$ is the current and $R$ is the resistance (\cref{ch4:sec:transmission}
shows where this comes from). A hotter wire has a higher resistance, so it heats even
more.

Engineers use \emph{thick copper wire}: a thick wire has less resistance, like a wide pipe
lets water through more easily. They use the thickest wire for the coil with the larger
current. They also \emph{cool} the transformer with air vents, fans or, in big
transformers, oil.

\subsection{Limitation 3: eddy currents in the core}

\term{Eddy currents} are loops of current induced inside a solid piece of metal when the
flux through it changes (\cref{ch3:sec:solenoids}). The iron core is a conductor, and its
flux changes all the time. So eddy currents circulate inside the core, in loops at right
angles to the flux (Faraday's law). By Lenz's law, their magnetic field opposes the change
in flux. They heat the core, and that energy never reaches the secondary. (In magnetic brakes, \cref{ch5:chap}, eddy currents are wanted. In a
transformer they only waste energy.)

The fix is a \term{laminated core}: a stack of thin iron sheets, each coated with a thin
insulating layer. The sheets lie \emph{along} the flux, so the flux passes easily. The
insulation cuts across the eddy-current paths, so each loop must stay inside one thin
sheet (\cref{ch4:fig:eddy}). Each thin loop surrounds only a thin strip of the changing
flux, so the emf that drives it is much smaller. The current must also squeeze through a
narrow path, which has a high resistance. So the eddy currents, and the heat they make,
are much smaller.

\begin{figure}
  \centering
  \begin{tikzpicture}[x=0.9cm, y=0.9cm]
    \tikzset{chFourdot/.pic={\draw[hscfield, line width=0.6pt, fill=white] (0,0) circle (0.12);
      \fill[hscfield] (0,0) circle (0.035);}}
    % (a) solid core
    \fill[chFouriron] (0,0) rectangle (3,3);
    \draw[hscink!70, line width=0.6pt] (0,0) rectangle (3,3);
    \draw[chFourcurrent, line width=1.3pt, rounded corners=8pt] (0.5,0.5) rectangle (2.5,2.5);
    \draw[chFourcurrent, line width=1.3pt, -{Stealth[length=2.6mm]}] (1.3,2.5) -- (1.75,2.5);
    \draw[chFourcurrent, line width=1.3pt, -{Stealth[length=2.6mm]}] (1.7,0.5) -- (1.25,0.5);
    \draw[chFourcurrent, line width=1.3pt, -{Stealth[length=2.6mm]}] (2.5,1.7) -- (2.5,1.25);
    \draw[chFourcurrent, line width=1.3pt, -{Stealth[length=2.6mm]}] (0.5,1.3) -- (0.5,1.75);
    \foreach \p in {(1.5,1.5),(0.25,0.25),(2.75,0.25),(0.25,2.75),(2.75,2.75)}
      \pic at \p {chFourdot};
    \node[font=\small\bfseries, text=hscink, align=center] at (1.5,3.65)
      {(a) solid core:\\ large eddy currents};
    % (b) laminated core
    \begin{scope}[xshift=4.8cm]
      \fill[chFouriron] (0,0) rectangle (3,3);
      \foreach \x in {0.5,1.0,...,2.6}
        \draw[hscink, line width=1.3pt] (\x,0) -- (\x,3);
      \draw[hscink!70, line width=0.6pt] (0,0) rectangle (3,3);
      \foreach \x in {0,0.5,...,2.6} {
        \draw[chFourcurrent, line width=1pt, rounded corners=3pt]
          (\x+0.13,0.35) rectangle (\x+0.37,2.65);
        \draw[chFourcurrent, line width=1pt, -{Stealth[length=2mm]}]
          (\x+0.37,1.7) -- (\x+0.37,1.3);
        \draw[chFourcurrent, line width=1pt, -{Stealth[length=2mm]}]
          (\x+0.13,1.3) -- (\x+0.13,1.7);
      }
      \node[font=\small\bfseries, text=hscink, align=center] at (1.5,3.65)
        {(b) laminated core:\\ small eddy currents};
      \node[diagramlabel, anchor=west] (ins) at (3.5,2.4) {insulating layer};
      \draw[hscmuted, line width=0.4pt] (ins.west) -- (2.5,2.2);
      \node[diagramlabel, anchor=west] (sh) at (3.5,1.5) {thin iron sheet};
      \draw[hscmuted, line width=0.4pt] (sh.west) -- (2.8,1.5);
      \node[diagramlabel, anchor=west] (ed) at (3.5,0.6) {eddy current};
      \draw[hscmuted, line width=0.4pt] (ed.west) -- (2.87,0.6);
      \foreach \p in {(0.25,0.15),(1.25,2.85),(2.75,0.15)}
        \pic at \p {chFourdot};
    \end{scope}
    \pic at (-0.45,-0.5) {chFourdot};
    \node[diagramlabel, anchor=west] at (-0.25,-0.5)
      {flux $\Phi$ points out of the page in both cores};
  \end{tikzpicture}
  \caption{A slice through the core. The flux comes out of the page towards you, and it is
    growing. The eddy currents flow clockwise (they reverse every half-cycle). Insulating
    layers cut each large loop into many thin ones.}
  \label{ch4:fig:eddy}
\end{figure}

Small transformers in chargers often use \term{ferrite} cores instead. Ferrite is a
magnetic ceramic that conducts electricity very poorly, so eddy currents are tiny.
Cooling does \emph{not} stop eddy currents; it only removes heat after it is made.

\textbf{Another loss: hysteresis.} Each cycle the core is magnetised one way, then the
other. This takes energy, called \term{hysteresis} loss, which also heats the core. Soft
iron is easy to magnetise and demagnetise, which keeps this loss low.

\textbf{Hum.} The core also vibrates a little. It does this twice in every cycle (once for
each direction of the field), so the hum is at $2 \times \qty{50}{Hz} = \qty{100}{Hz}$.

\textbf{Is the ideal model good enough?} The \term{efficiency} ($\eta$) is the percentage
of the input power that comes out as useful power. Big transformers in power stations
and substations (places where the grid changes voltage, \cref{ch4:sec:transmission}) are
about 99\% efficient. Small ones often waste 10--20\%. So the ideal model is a very good
approximation, but never exact. Engineers can reduce the losses, but never remove them.

\begin{toolkit}{Percentages and decimals}
To change a percentage to a decimal, divide by 100: $92\% = \dfrac{92}{100} = 0.92$. To
find 85\% of \qty{200}{W}: $0.85 \times \qty{200}{W} = \qty{170}{W}$. To change a decimal
back to a percentage, multiply by 100.
\end{toolkit}

\begin{formulacard}[notonsheet]{Efficiency}
  \cardformula{\eta = \frac{P_\mathrm{out}}{P_\mathrm{in}} \times 100\%
    = \frac{V_\mathrm{s}I_\mathrm{s}}{V_\mathrm{p}I_\mathrm{p}} \times 100\%}
  \begin{symbols}
    \sym{$\eta$}{efficiency (Greek letter ``eta'')}{no unit (\%)}
    \sym{$P_\mathrm{out}$, $P_\mathrm{in}$}{useful power out of the secondary, power into
      the primary}{\si{W}}
  \end{symbols}
  \inwords{the percentage of the input power that comes out as useful power.}
  \whentouse{a question says a transformer loses energy, or gives its efficiency.}
  \whennottouse{an ideal transformer: then $\eta = 100\%$.}
  \rearranged{$P_\mathrm{out} = \eta P_\mathrm{in}$ \qquad
    $I_\mathrm{s} = \dfrac{\eta V_\mathrm{p}I_\mathrm{p}}{V_\mathrm{s}}$ \qquad
    ($\eta$ as a decimal, for example $0.92$)}
  \tensecond{\qty{200}{W} in and \qty{190}{W} out, so
    $\eta = \dfrac{190}{200} \times 100\% = 95\%$.}
  \derive{it is a definition: useful power out $\div$ power in. The power lost is
    $P_\mathrm{in} - P_\mathrm{out}$.}
  \trap{the efficiency goes on the \emph{power}, not on the voltage. Use $0.92$, not $92$,
    inside a formula.}
\end{formulacard}

\begin{formulacard}{Power and energy}
  \cardformula{P = \frac{\Delta E}{\Delta t}}
  \begin{symbols}
    \sym{$P$}{power}{\si{W}}
    \sym{$\Delta E$}{energy transferred (for example, wasted as heat)}{\si{J}}
    \sym{$\Delta t$}{time taken}{\si{s}}
  \end{symbols}
  \inwords{power is the energy transferred each second.}
  \whentouse{to find the energy used or wasted in a given time, or the power from an
    energy and a time.}
  \whennottouse{with a time in minutes or hours. Change it to seconds first.}
  \rearranged{$\Delta E = P\,\Delta t$ \qquad $\Delta t = \dfrac{\Delta E}{P}$}
  \tensecond{a transformer wastes \qty{25}{W} for \qty{2.0}{min} $= \qty{120}{s}$, so
    $\Delta E = 25 \times 120 = \qty{3000}{J}$.}
  \trap{change hours or minutes to seconds: $\qty{1}{h} = 60 \times 60 = \qty{3600}{s}$.
    Use the power \emph{wasted} ($P_\mathrm{in} - P_\mathrm{out}$) for the energy wasted.}
\end{formulacard}

\begin{workedexample}{How much energy is wasted?}\label{ch4:we:eff}
A transformer takes in \qty{2400}{W} and gives out \qty{2280}{W}. (a)~Calculate its
efficiency. (b)~How much energy is wasted in one hour?
\given{$P_\mathrm{in} = \qty{2400}{W}$, \ $P_\mathrm{out} = \qty{2280}{W}$, \
  $\Delta t = \qty{1}{h} = \qty{3600}{s}$}
\find{(a) $\eta$ \quad (b) the energy wasted, $\Delta E$}
\formula{$\eta = \dfrac{P_\mathrm{out}}{P_\mathrm{in}} \times 100\%$ \quad and \quad
  $P = \dfrac{\Delta E}{\Delta t}$, so $\Delta E = P\,\Delta t$}
\substitute{(a) $\eta = \dfrac{\qty{2280}{W}}{\qty{2400}{W}} \times 100\% = 95.0\%$\par
  (b) power wasted $= \qty{2400}{W} - \qty{2280}{W} = \qty{120}{W}$, so
  $\Delta E = \qty{120}{W} \times \qty{3600}{s} = \qty{432000}{J}$}
\answer{(a) \final{$\eta = 95.0\%$} \quad (b) \final{$\Delta E = \qty{4.32e5}{J}$}}
\checkstep{$\eta < 100\%$ \cmark. Wasting \qty{120}{W} is like leaving a bright old light
  globe on, which is enough to make the transformer warm.}
\end{workedexample}

\begin{workedexample}{A real step-down transformer}\label{ch4:we:real}
A transformer for \qty{12}{V} downlights in a kitchen is 92\% efficient. Its primary is
at \qty{230}{V} and carries \qty{0.50}{A}. The output voltage is \qty{12}{V}. (a)~Calculate the secondary current.
(b)~Compare it with the ideal value and give the main cause of the difference.
\given{$\eta = 0.92$, \ $V_\mathrm{p} = \qty{230}{V}$, \ $I_\mathrm{p} = \qty{0.50}{A}$, \
  $V_\mathrm{s} = \qty{12}{V}$}
\find{(a) $I_\mathrm{s}$ \quad (b) the ideal $I_\mathrm{s}$, and why they differ}
\formula{$P_\mathrm{in} = V_\mathrm{p}I_\mathrm{p}$, \ $P_\mathrm{out} = \eta P_\mathrm{in}$,
  \ $I_\mathrm{s} = \dfrac{P_\mathrm{out}}{V_\mathrm{s}}$}
\substitute{(a) $P_\mathrm{in} = \qty{230}{V} \times \qty{0.50}{A} = \qty{115}{W}$, \
  $P_\mathrm{out} = 0.92 \times \qty{115}{W} = \qty{105.8}{W}$, \
  $I_\mathrm{s} = \dfrac{\qty{105.8}{W}}{\qty{12}{V}} = \qty{8.82}{A}$\par
  (b) ideal: $I_\mathrm{s} = \dfrac{\qty{115}{W}}{\qty{12}{V}} = \qty{9.58}{A}$}
\answer{(a) \final{$I_\mathrm{s} = \qty{8.8}{A}$} \quad (b) ideal \qty{9.6}{A}. The
  missing \qty{9.2}{W} is mostly heat from resistance in the coils and eddy currents in
  the core.}
\checkstep{$\dfrac{8.82}{9.58} = 0.92$, which equals $\eta$ \cmark. The real current is a
  little smaller than the ideal one \cmark. Step-down, so $I_\mathrm{s} > I_\mathrm{p}$
  \cmark.}
\end{workedexample}

\begin{mistakes}
\begin{itemize}
  \item \xmark\ ``Energy is destroyed.'' \cmark\ It becomes heat and sound; energy is
        conserved.
  \item \xmark\ ``Cooling reduces eddy currents.'' \cmark\ Laminations (or ferrite) do.
        Cooling only removes heat.
  \item \xmark\ Laminations across the flux. \cmark\ The sheets lie along the flux; the
        insulation blocks the eddy currents.
  \item \xmark\ Naming a loss without the physics. \cmark\ Give the cause, the effect and
        the fix.
\end{itemize}
\end{mistakes}

\begin{examtip}
``Evaluate the ideal transformer model'' needs a judgement. (1)~State that the model
assumes no losses ($V_\mathrm{p}I_\mathrm{p} = V_\mathrm{s}I_\mathrm{s}$). (2)~Explain two
or three limitations with their causes: flux leakage, $I^2R$ heating in the coils, eddy
currents in the core. (3)~Give a fix for each: closed iron core, thick copper wire,
laminated core. (4)~Judge: ``A very good approximation, but real output power is always a
little less than input power.''
\end{examtip}

\begin{quickcheck}
  \item A transformer takes in \qty{500}{W} and gives out \qty{470}{W}. Find its
        efficiency, and the energy it wastes in 10 minutes.
  \item Explain why the core of a transformer is laminated instead of
        solid.\label{ch4:qc:laminate}
  \item A student puts two coils side by side on a bench, with no iron core. The turns
        equation predicts a secondary voltage of \qty{6.0}{V}, but the voltmeter reads only
        \qty{1.5}{V}. Name the limitation that causes this, explain it, and give one way to
        fix it.
  \item A transformer is 85\% efficient. $V_\mathrm{p} = \qty{230}{V}$,
        $I_\mathrm{p} = \qty{2.0}{A}$ and $V_\mathrm{s} = \qty{46}{V}$. Find the
        secondary current.
  \item A wireless phone charger has an air gap and no shared iron core. Explain why it
        wastes more energy than a transformer with a closed core.\label{ch4:qc:wireless}
\end{quickcheck}

%% ==================================================================
\section{Sending electricity at high voltage}\label{ch4:sec:transmission}
\covers{M6.IQ3.05}

\begin{bigidea}
Power lines have resistance, so the current in them makes heat. This loss is $I^2R$: it
depends on the current squared. Stepping the voltage up makes the current small, so much
less energy is wasted.
\end{bigidea}

\begin{keywords}
  \kw{Transmission line}{a long, thick wire on tall towers that carries electricity far.}
  \kw{Voltage drop}{the voltage used up pushing current through a resistance, $IR$.}
  \kw{Substation}{a place with transformers that change the voltage (near towns, they
    step it down).}
  \kw{Electricity grid}{the network of power stations, lines and transformers.}
\end{keywords}

Most power stations and wind farms are far from cities. \term{Transmission lines}, long
thick wires on tall towers, carry the electricity. Even good wires have some resistance,
and a line hundreds of kilometres long can have several ohms.

\Cref{ch4:fig:line} models the line as one resistance $R$. Some voltage is used up
pushing the current $I$ through this resistance. This is the \term{voltage drop}:
$V_\mathrm{drop} = IR$ (Ohm's law). The power turned into heat in the line is the power
loss, $P_\mathrm{loss}$. Less current means much less heat, as the formula card below
shows.

\begin{figure}
  \centering
  \begin{tikzpicture}[x=1cm, y=1cm]
    \draw[hscink, line width=0.7pt] (0,0.95) -- (0,1.6) -- (2.6,1.6);
    \draw[hscink, line width=0.7pt] (4.4,1.6) -- (7.4,1.6) -- (7.4,1.1);
    \draw[hscink, line width=0.7pt] (7.4,0.1) -- (7.4,-0.4) -- (0,-0.4) -- (0,0.25);
    \draw[hscink, line width=0.7pt, fill=white] (0,0.6) circle (0.35);
    \node at (0,0.6) {$\sim$};
    \draw[hscink, line width=0.7pt, fill=hscforce!10] (2.6,1.4) rectangle (4.4,1.8);
    \node[diagramlabel] at (3.5,1.6) {$R$};
    \draw[hscink, line width=0.7pt, fill=hscbrand!8] (7.0,0.1) rectangle (7.8,1.1);
    \node[diagramlabel, rotate=90] at (7.4,0.6) {town};
    \draw[velocity] (1.0,1.6) -- (1.9,1.6) node[midway, above, text=hscvel] {$I$};
    \node[diagramlabel, align=center, anchor=east] at (-0.45,0.6)
      {power station\\ sends voltage $V$};
    \node[diagramlabel, align=center, above] at (3.5,1.85)
      {line resistance $R$ (both wires)\\ voltage drop $V_\mathrm{drop} = IR$};
    \node[diagramlabel, align=center, anchor=west] at (7.95,0.6)
      {voltage left\\ for the town\\ $= V - IR$};
    \node[diagramlabel, text=hscforce] at (3.5,1.0) {heat: $P_\mathrm{loss} = I^2R$};
  \end{tikzpicture}
  \caption{A simple model of a transmission line. The current $I$ is the same all the way
    around the loop.}
  \label{ch4:fig:line}
\end{figure}

\begin{toolkit}{Prefixes, scientific notation, squares and square roots}
k (kilo) $= 10^3$ and M (mega) $= 10^6$. Change to the base unit before you calculate:
$\qty{330}{kV} = 330 \times 10^3\ \si{V}$. Move the decimal point 2 places to the left and
add 2 to the power: $\qty{3.30e5}{V}$. In the same way,
$\qty{20}{MW} = 20 \times 10^6\ \si{W} = \qty{2.0e7}{W}$.\par\smallskip
Squaring a change: if the current is multiplied by a number, the loss $I^2R$ is multiplied
by that number squared. So $\tfrac{1}{4}$ of the current gives
$\left(\tfrac{1}{4}\right)^2 = \tfrac{1}{16}$ of the loss. To undo a square, take the
square root of both sides: if $I^2 = 900$, then $I = \sqrt{900} = 30$.
\end{toolkit}

\begin{formulacard}[notonsheet]{Power lost in a transmission line}
  \cardformula{P_\mathrm{loss} = I^2R}
  \begin{symbols}
    \sym{$P_\mathrm{loss}$}{power turned into heat in the line}{\si{W}}
    \sym{$I$}{current in the line}{\si{A}}
    \sym{$R$}{total resistance of the line (both wires)}{\si{\ohm}}
  \end{symbols}
  \inwords{the heat lost each second grows with the \emph{square} of the current. Halve
    the current and the loss falls to a quarter.}
  \whentouse{heat lost in power lines, extension leads or transformer coils. Find the
    current first with $I = \dfrac{P}{V}$.}
  \whennottouse{for the power the town receives. That is $P - P_\mathrm{loss}$.}
  \rearranged{$I = \sqrt{\dfrac{P_\mathrm{loss}}{R}}$ \qquad
    $R = \dfrac{P_\mathrm{loss}}{I^2}$}
  \tensecond{$I = \qty{100}{A}$ in a \qty{3.0}{\ohm} line, so
    $P_\mathrm{loss} = 100^2 \times 3.0 = \qty{3.0e4}{W}$.}
  \derive{combine $P = VI$ and $V = IR$ (both on the sheet), using the voltage
    \emph{drop} across the line.
    \begin{align*}
      P_\mathrm{loss} &= V_\mathrm{drop}\, I && \text{$P = VI$ for the line only}\\
      &= (IR)\, I && \text{the drop across the line is $V_\mathrm{drop} = IR$}\\
      &= I^2R && \text{multiply the two $I$s}
    \end{align*}}
  \trap{do not use $\dfrac{V^2}{R}$ with the line's voltage (for example \qty{330}{kV}).
    There, $V$ must be the voltage drop $IR$, not the voltage sent.}
\end{formulacard}

\textbf{Why high voltage?} A power station sends a fixed power $P$. Since $P = VI$, a
higher voltage needs a smaller current. The loss $I^2R$ depends on the \emph{square} of
the current. Example: 10 times the voltage gives $\tfrac{1}{10}$ of the current. Then the
loss is $\left(\tfrac{1}{10}\right)^2 = \tfrac{1}{100}$ of what it was. Thicker wires
would also help, but they are heavy and very expensive.

\textbf{The grid.} The \term{electricity grid} is the network of power stations, lines
and transformers (\cref{ch4:fig:grid}). A \term{substation} is a place with transformers
that change the voltage (near towns, they step it down).

\begin{sidebyside}[0.45]
  \centering
  \begin{tikzpicture}[x=1cm, y=1cm, node distance=3.8mm,
      stage/.style={concept=33mm, font=\footnotesize, inner sep=1.1mm},
      tlabel/.style={font=\footnotesize, text=hscbrand, anchor=west, align=left}]
    \node[stage] (a) {power station\\ generator: about \qty{20}{kV}};
    \node[stage, below=of a] (b) {transmission lines\\ \qty{330}{kV} or \qty{500}{kV}};
    \node[stage, below=of b] (c) {lines to suburbs\\ \qty{132}{kV} or \qty{66}{kV}};
    \node[stage, below=of c] (d) {local lines \qty{11}{kV}};
    \node[stage, below=of d] (e) {homes \qty{230}{V}};
    \draw[link] (a) -- node[tlabel, right=1mm] {step-up} (b);
    \draw[link] (b) -- node[tlabel, right=1mm] {step-down (substation)} (c);
    \draw[link] (c) -- node[tlabel, right=1mm] {step-down (substation)} (d);
    \draw[link] (d) -- node[tlabel, right=1mm] {step-down (on a power pole)} (e);
  \end{tikzpicture}
  \captionof{figure}{From power station to home (typical NSW voltages). The last
    transformer sits on a power pole in your street.}
  \label{ch4:fig:grid}
\nextside
  \begin{itemize}
    \item \textbf{Step up} at the power station: a low current in the long lines, so
          little energy is lost.
    \item \textbf{Step down} near towns: high voltages are dangerous. They can spark
          (arc) across air, and need tall towers and big insulators.
    \item Homes get \qty{230}{V}, the voltage appliances are built for.
    \item \textbf{Why AC?} Transformers need a changing current, so they work with AC
          but not with steady DC. (A few undersea links, such as Basslink, use DC.)
  \end{itemize}
\end{sidebyside}

\textbf{Everyday uses.} Transformers are not only used to send electricity far.
\Cref{ch4:tab:uses} shows some step-up and step-down transformers you may meet.

\begin{table}
  \caption{Some uses of step-up and step-down transformers.}\label{ch4:tab:uses}
  \centering\small
  \begin{tabular}{@{}L{27mm}L{17mm}L{26mm}L{65mm}@{}}
    \toprule
    \colhead{Where} & \colhead{Type} & \colhead{About} & \colhead{Why}\\
    \midrule
    Power station & step-up & \qty{20}{kV} $\to$ \qty{330}{kV} & less loss in long lines\\
    Microwave oven & step-up & \qty{230}{V} $\to$ \qty{2}{kV} & high voltage makes the microwaves\\
    Welding machine & step-down & \qty{230}{V} $\to$ \qty{50}{V} & very large current melts metal\\
    Phone charger & step-down & \qty{230}{V} $\to$ \qty{5}{V} & low voltage for the phone: a
      small high-frequency ferrite-core transformer, then electronics make DC\\
    \bottomrule
  \end{tabular}
\end{table}

\begin{workedexample}{Power to a farm}\label{ch4:we:farm}
A small power station sends \qty{60.0}{kW} to a farm at \qty{2.40}{kV}. The line has a
total resistance of \qty{4.0}{\ohm}. Calculate (a)~the current, (b)~the power lost in the
line, (c)~the power that reaches the farm and (d)~the voltage that reaches the farm.
(e)~At the farm, an ideal step-down transformer gives \qty{230}{V}. Find its turns ratio.
\given{$P = \qty{6.00e4}{W}$, \ $V = \qty{2.40}{kV} = \qty{2400}{V}$, \
  $R = \qty{4.0}{\ohm}$, \ farm voltage $V_\mathrm{s} = \qty{230}{V}$}
\find{(a) $I$ \quad (b) $P_\mathrm{loss}$ \quad (c) the power received \quad
  (d) $V_\mathrm{farm}$ \quad (e) $N_\mathrm{p} : N_\mathrm{s}$}
\formula{$I = \dfrac{P}{V}$, \ then $P_\mathrm{loss} = I^2R$, \ then
  received $= P - P_\mathrm{loss}$. \ For (d): $V_\mathrm{drop} = IR$ and
  $V_\mathrm{farm} = V - V_\mathrm{drop}$. \ For (e):
  $\dfrac{N_\mathrm{p}}{N_\mathrm{s}} = \dfrac{V_\mathrm{p}}{V_\mathrm{s}}$, with
  $V_\mathrm{p} = V_\mathrm{farm}$.}
\substitute{(a) $I = \dfrac{\qty{6.00e4}{W}}{\qty{2400}{V}} = \qty{25.0}{A}$\par
  (b) $P_\mathrm{loss} = (\qty{25.0}{A})^2 \times \qty{4.0}{\ohm} = \qty{2500}{W}$\par
  (c) received $= \qty{60.0}{kW} - \qty{2.5}{kW} = \qty{57.5}{kW}$\par
  (d) $V_\mathrm{drop} = \qty{25.0}{A} \times \qty{4.0}{\ohm} = \qty{100}{V}$, so
  $V_\mathrm{farm} = \qty{2400}{V} - \qty{100}{V} = \qty{2300}{V}$\par
  (e) $\dfrac{N_\mathrm{p}}{N_\mathrm{s}} = \dfrac{\qty{2300}{V}}{\qty{230}{V}} = 10$}
\answer{(a) \final{$I = \qty{25.0}{A}$} \quad (b) \final{$P_\mathrm{loss} = \qty{2.5}{kW}$}
  \quad (c) \final{\qty{57.5}{kW}} \quad (d) \final{\qty{2300}{V}} \quad
  (e) \final{$N_\mathrm{p} : N_\mathrm{s} = 10 : 1$}}
\checkstep{the loss is $\dfrac{2.5}{60.0} = 4.2\%$ of the power sent: sensible \cmark.
  The wrong method, $\dfrac{V^2}{R} = \dfrac{2400^2}{4.0} = \qty{1.44}{MW}$, is more than
  the power sent: impossible.\par
  The current on the farm side is
  $\dfrac{\qty{57500}{W}}{\qty{230}{V}} = \qty{250}{A} = 10 \times \qty{25.0}{A}$
  \cmark. Step-down, so the current goes up by the turns ratio.}
\end{workedexample}

\begin{workedexample}{Why a wind farm steps up the voltage}\label{ch4:we:wind}
A wind farm sends \qty{20.0}{MW} along a line of total resistance \qty{8.0}{\ohm}.
(a)~Find the power lost if it sends at \qty{22.0}{kV}. (b)~A transformer with turns ratio
$N_\mathrm{p} : N_\mathrm{s} = 1 : 6$ first steps the voltage up. Find the new voltage
and the new loss. (c)~By what factor does the loss fall?
\given{$P = \qty{2.00e7}{W}$, \ $R = \qty{8.0}{\ohm}$, \ $V_1 = \qty{2.20e4}{V}$, \
  $N_\mathrm{p} : N_\mathrm{s} = 1 : 6$}
\find{(a) $P_\mathrm{loss}$ at \qty{22.0}{kV} \quad (b) $V_2$ and the new $P_\mathrm{loss}$
  \quad (c) the factor}
\formula{$I = \dfrac{P}{V}$, \ $P_\mathrm{loss} = I^2R$, \
  $V_\mathrm{s} = V_\mathrm{p}\dfrac{N_\mathrm{s}}{N_\mathrm{p}}$}
\substitute{(a) $I_1 = \dfrac{\qty{2.00e7}{W}}{\qty{2.20e4}{V}} = \qty{909}{A}$, \
  $P_\mathrm{loss} = (\qty{909}{A})^2 \times \qty{8.0}{\ohm} = \qty{6.61e6}{W}$\par
  (b) $V_2 = \qty{22.0}{kV} \times \dfrac{6}{1} = \qty{132}{kV}$, \
  $I_2 = \dfrac{\qty{2.00e7}{W}}{\qty{1.32e5}{V}} = \qty{151.5}{A}$, \
  $P_\mathrm{loss} = (\qty{151.5}{A})^2 \times \qty{8.0}{\ohm} = \qty{1.84e5}{W}$\par
  (c) $\dfrac{\qty{6.61e6}{W}}{\qty{1.84e5}{W}} = 35.9 \approx 36$}
\answer{(a) \final{\qty{6.6}{MW}} (33\%) \quad
  (b) \final{\qty{132}{kV}}, \final{\qty{0.18}{MW}} (about 0.9\%)\par
  (c) \final{36 times}}
\checkstep{6 times the voltage gives $\tfrac{1}{6}$ of the current, so the loss falls by
  $6^2 = 36$ \cmark.}
\end{workedexample}

\begin{mistakes}
\begin{itemize}
  \item \xmark\ Using $\dfrac{V^2}{R}$ with the transmission voltage.
        \cmark\ Find $I = \dfrac{P}{V}$ first, then use $I^2R$.
  \item \xmark\ ``The loss is proportional to $I$.'' \cmark\ It is proportional to $I^2$.
  \item \xmark\ ``High-voltage lines have less resistance.'' \cmark\ Same $R$, smaller $I$.
\end{itemize}
\end{mistakes}

\begin{examtip}
``Justify the use of high-voltage AC transmission'' (4 marks). Link every claim to its
reason, and use numbers when you can.
\begin{enumerate}
  \item For a fixed power, $P = VI$. So a higher voltage means a smaller current.
  \item $P_\mathrm{loss} = I^2R$. So the loss falls with the square of the current.
  \item Transformers can step AC up and down, because AC makes a changing flux. They
        cannot do this with steady DC.
  \item The voltage is stepped down near homes for safety.
\end{enumerate}
\end{examtip}

\begin{quickcheck}
  \item A line of resistance \qty{2.0}{\ohm} carries \qty{150}{A}. Find the power lost.
  \item A power station sends the same power at 4 times the voltage. What happens to the
        current and to the power lost in the line?
  \item A \qty{6.0}{\ohm} line carries \qty{200}{MW} at \qty{330}{kV}. A student finds
        the loss as $\dfrac{(\num{330000})^2}{6.0} = \qty{1.8e10}{W}$. Explain the error,
        and find the correct loss.\label{ch4:qc:v2r}
  \item Explain why the grid uses AC, and why the voltage is stepped down before it
        reaches homes.
  \item A power station sends \qty{200}{A} at \qty{67.0}{kV} along a line of total
        resistance \qty{5.0}{\ohm}. (a)~Find the voltage at the far end. (b)~A substation
        transformer there gives \qty{11.0}{kV}. Is it step-up or step-down? Find
        $N_\mathrm{p} : N_\mathrm{s}$. (c)~Assume it is ideal. Find the current in the
        \qty{11.0}{kV} lines.\label{ch4:qc:substation}
\end{quickcheck}

%% ==================================================================
\section{Key experiment}\label{ch4:sec:experiments}
\covers{M6.IQ3.03, M6.IQ3.04}

\begin{toolkit}{Gradient from a graph}
The gradient of a straight line is how steep it is: how much $y$ goes up for each 1 unit
along $x$. Pick two points on the line, far apart. Then
gradient $= \dfrac{\text{rise}}{\text{run}} = \dfrac{y_2 - y_1}{x_2 - x_1}$. Example
(\cref{ch4:fig:vsns}): the line of best fit passes through $(0, \qty{0}{V})$ and
$(1200, \qty{11.4}{V})$, so gradient
$= \dfrac{(11.4 - 0)\ \si{V}}{(1200 - 0)\ \text{turns}} = \qty{0.0095}{V}$ per turn.
A straight line through the origin means $y$ is \term{directly proportional} to $x$,
written $y \propto x$: double $x$ and $y$ doubles.
\end{toolkit}

\begin{experiment}{Testing the transformer equations}
\aim{To test $V_\mathrm{p}/V_\mathrm{s} = N_\mathrm{p}/N_\mathrm{s}$ for a real
  transformer, and to find its efficiency.}
\apparatus{a transformer kit that comes apart (a \emph{demountable} kit): a laminated
  U-shaped iron core, a straight iron bar (the \emph{I-bar}) that closes the top of the U,
  and coils of 150 to 1200 turns that slide onto the two arms (sides) of the U; a
  low-voltage AC power pack; two AC voltmeters and two AC ammeters (meters made to read
  AC); a load resistor (\cref{ch4:fig:apparatus}).
  \begin{figure}
    \centering
    \begin{tikzpicture}[x=0.9cm, y=0.9cm]
      % U-core
      \fill[chFouriron] (0,0) -- (3.2,0) -- (3.2,2.3) -- (2.5,2.3) -- (2.5,0.7) --
        (0.7,0.7) -- (0.7,2.3) -- (0,2.3) -- cycle;
      \draw[hscink!70, line width=0.6pt] (0,0) -- (3.2,0) -- (3.2,2.3) -- (2.5,2.3) --
        (2.5,0.7) -- (0.7,0.7) -- (0.7,2.3) -- (0,2.3) -- cycle;
      % I-bar, lifted above the core
      \fill[chFouriron] (0,2.75) rectangle (3.2,3.4);
      \draw[hscink!70, line width=0.6pt] (0,2.75) rectangle (3.2,3.4);
      \draw[displacement] (1.6,2.7) -- (1.6,2.38);
      % coils on the two arms
      \foreach \cx in {0.35,2.85} {
        \fill[white] (\cx-0.5,0.95) rectangle (\cx+0.5,2.05);
        \draw[chFourcopper, line width=0.8pt] (\cx-0.5,0.95) rectangle (\cx+0.5,2.05);
        \foreach \y in {1.05,1.17,...,1.98}
          \draw[chFourcopper, line width=0.6pt] (\cx-0.5,\y) -- (\cx+0.5,\y);
      }
      % labels
      \node[diagramlabel, anchor=west] (ib) at (3.4,3.07) {I-bar};
      \node[diagramlabel, anchor=north] at (1.6,-0.08) {U-shaped core};
      \node[diagramlabel, anchor=east, align=right] at (-0.3,1.5) {primary\\ coil};
      \node[diagramlabel, anchor=west, align=left] at (3.4,1.5) {secondary\\ coil};
      \node[diagramlabel, anchor=west] (arm) at (3.4,0.35) {arm};
      \draw[hscmuted, line width=0.4pt] (arm.west) -- (2.95,0.55);
      \node[font=\small\bfseries, text=hscink] at (1.6,4.0) {(a) the kit};
    \end{tikzpicture}\hspace{4mm}%
    \begin{tikzpicture}[x=0.8cm, y=0.8cm, line width=0.7pt]
      \tikzset{meter/.style={draw=hscink, fill=white, circle, minimum size=6.5mm,
        inner sep=0pt, font=\small\bfseries}}
      % primary circuit
      \draw (0,0.3) -- (0,0) -- (3,0) -- (3,0.4);
      \draw (0,0.9) -- (0,2) -- (3,2) -- (3,1.6);
      \draw (3,1.6) arc[start angle=90, end angle=-90, radius=0.15]
        arc[start angle=90, end angle=-90, radius=0.15]
        arc[start angle=90, end angle=-90, radius=0.15]
        arc[start angle=90, end angle=-90, radius=0.15];
      \draw[fill=white] (0,0.6) circle (0.3); \node at (0,0.6) {$\sim$};
      \node[meter] at (0.75,2) {A};
      \draw (1.65,2) -- (1.65,0);
      \node[meter] at (1.65,1) {V};
      % core
      \draw[line width=1.2pt] (3.42,0.3) -- (3.42,1.7) (3.58,0.3) -- (3.58,1.7);
      % secondary circuit
      \draw (4.0,1.6) arc[start angle=90, end angle=270, radius=0.15]
        arc[start angle=90, end angle=270, radius=0.15]
        arc[start angle=90, end angle=270, radius=0.15]
        arc[start angle=90, end angle=270, radius=0.15];
      \draw (4.0,1.6) -- (4.0,2) -- (6.6,2) -- (6.6,1.45);
      \draw (4.0,0.4) -- (4.0,0) -- (6.6,0) -- (6.6,0.55);
      \draw[fill=hscbrand!8] (6.4,0.55) rectangle (6.8,1.45);
      \draw (5.3,2) -- (5.3,0);
      \node[meter] at (5.3,1) {V};
      \node[meter] at (6.05,2) {A};
      % labels
      \node[diagramlabel, align=center, anchor=north] at (0,-0.1) {AC power\\ pack};
      \node[diagramlabel, anchor=east] at (2.75,1.0) {$N_\mathrm{p}$};
      \node[diagramlabel, anchor=west] at (4.2,1.0) {$N_\mathrm{s}$};
      \node[diagramlabel] at (3.5,2.35) {iron core};
      \node[diagramlabel, align=center, anchor=north] at (6.6,-0.1) {load\\ resistor};
      \node[font=\small\bfseries, text=hscink] at (3.3,3.2) {(b) the circuit};
    \end{tikzpicture}
    \caption{(a) The demountable kit: two coils slide onto the arms of the U-core, and
      the I-bar closes the top. (b) The test circuit. The symbol in the middle is an
      iron-core transformer.}
    \label{ch4:fig:apparatus}
  \end{figure}}
\method{\begin{enumerate}
  \item Slide a 600-turn primary onto one arm of the U-core and join it to \qty{6.0}{V}
        AC. Slide a secondary coil onto the other arm. Put the I-bar across the top to
        close the core.
  \item With no load (only the voltmeter joined to the secondary), read $V_\mathrm{s}$
        for secondaries of 150 to 1200 turns, 3 times each. Switch off before you change
        a coil.
  \item With the load on a 300-turn secondary, read all four meters ($V_\mathrm{p}$,
        $I_\mathrm{p}$, $V_\mathrm{s}$, $I_\mathrm{s}$). Then
        $P_\mathrm{in} = V_\mathrm{p}I_\mathrm{p}$ and
        $P_\mathrm{out} = V_\mathrm{s}I_\mathrm{s}$.
  \item Switch off, remove the load and lift off the I-bar. Switch on and read
        $V_\mathrm{s}$ again.
\end{enumerate}}
\results{Sample results from a school kit:
  \begin{center}\small\vspace{-6pt}
  \begin{tabular}{@{}lccccc@{}}
    \toprule
    $N_\mathrm{s}$ (with $N_\mathrm{p} = 600$, $V_\mathrm{p} = \qty{6.0}{V}$)
      & 150 & 300 & 600 & 900 & 1200\\
    \midrule
    ideal $V_\mathrm{s} = V_\mathrm{p}N_\mathrm{s}/N_\mathrm{p}$ (\si{V})
      & 1.50 & 3.00 & 6.00 & 9.00 & 12.00\\
    measured $V_\mathrm{s}$, mean of 3 (\si{V}) & 1.42 & 2.86 & 5.71 & 8.55 & 11.36\\
    \bottomrule
  \end{tabular}
  \end{center}
  \begin{figure}
    \centering
    \begin{tikzpicture}
      \begin{axis}[hsc, width=86mm, height=44mm, xmin=0, xmax=1300, ymin=0, ymax=13.5,
          xtick={0,300,600,900,1200}, ytick={0,3,6,9,12},
          xticklabel style={/pgf/number format/1000 sep={}},
          xlabel={Number of secondary turns $N_\mathrm{s}$},
          ylabel={$V_\mathrm{s}$ (\si{V})}]
        \addplot[hscseriesB, dashed, line width=1pt, domain=0:1300, samples=2] {0.010*x};
        \addplot[hscseriesA, line width=1pt, domain=0:1300, samples=2] {0.009485*x};
        \addplot[hscseriesA, only marks, mark=*, mark size=2pt]
          coordinates {(150,1.42) (300,2.86) (600,5.71) (900,8.55) (1200,11.36)};
        \node[font=\small, text=hscseriesB, anchor=south east] at (axis cs:1150,12.3)
          {ideal: \qty{0.010}{V} per turn};
        \node[font=\small, text=hscseriesA, anchor=north west] at (axis cs:1000,9.0)
          {best fit: \qty{0.0095}{V} per turn};
      \end{axis}
    \end{tikzpicture}
    \caption{Measured points lie just below the ideal line (dashed).}
    \label{ch4:fig:vsns}
  \end{figure}
  Load test (300-turn secondary): $V_\mathrm{p} = \qty{6.0}{V}$,
  $I_\mathrm{p} = \qty{0.52}{A}$, $V_\mathrm{s} = \qty{2.80}{V}$,
  $I_\mathrm{s} = \qty{1.00}{A}$. So
  $P_\mathrm{in} = \qty{6.0}{V} \times \qty{0.52}{A} = \qty{3.12}{W}$ (\qty{3.1}{W} to 2
  significant figures), $P_\mathrm{out} = \qty{2.80}{V} \times \qty{1.00}{A}
  = \qty{2.80}{W}$ and
  $\eta = \dfrac{\qty{2.80}{W}}{\qty{3.12}{W}} \times 100\% = 89.7\% \approx 90\%$.
  Open core (no load): $V_\mathrm{s}$ fell from \qty{2.86}{V} to about \qty{0.9}{V}.}
\conclusion{$V_\mathrm{s}$ against $N_\mathrm{s}$ is a straight line through the origin
  ($V_\mathrm{s} \propto N_\mathrm{s}$), as the turns equation predicts. Its gradient,
  \qty{0.0095}{V} per turn, is 5\% below the ideal
  $V_\mathrm{p}/N_\mathrm{p} = \qty{0.010}{V}$ per turn, mainly because some flux leaks
  (incomplete flux linkage). With a load, $P_\mathrm{out} < P_\mathrm{in}$
  ($\eta = 90\%$): energy is lost as heat. The ideal model is close, but not exact.}
\significance{It supports the model behind every grid transformer, and shows why cores
  are closed loops. (Faraday's iron ring, the first transformer, is in
  \cref{ch3:sec:experiments}.)}
\validity{\begin{itemize}
  \item \textbf{Validity:} change only $N_\mathrm{s}$; keep $V_\mathrm{p}$, $N_\mathrm{p}$,
        the frequency and the core the same; use AC meters.
  \item \textbf{Reliability:} repeat each reading 3 times and average; use 5 coils.
  \item \textbf{Accuracy:} use calibrated AC meters on a suitable range. Every measured
        $V_\mathrm{s}$ is about 5\% below the ideal value. A gap that is the same for every
        coil is a \emph{systematic} error (here, mainly flux leakage); repeating cannot
        remove it. \emph{Random} errors make readings scatter; averaging reduces them.
  \item \textbf{Safety:} never use mains; a $1:10$ step-up turns \qty{12}{V} into
        \qty{120}{V}. Switch off the power pack before you change a coil or move the
        I-bar.
\end{itemize}}
\end{experiment}

%% ==================================================================
\begin{summarybox}{\Cref{ch4:chap} in five lines}
\begin{itemize}
  \item A transformer needs AC. A changing current in the primary makes a changing flux
        in the iron core. This induces an emf in the secondary (Faraday's law).
  \item Ideal transformer: $V_\mathrm{p}/V_\mathrm{s} = N_\mathrm{p}/N_\mathrm{s}$ and
        $V_\mathrm{p}I_\mathrm{p} = V_\mathrm{s}I_\mathrm{s}$. Voltage up means current
        down; the frequency stays the same.
  \item Real transformers are not ideal: flux leakage (some flux misses the secondary),
        $I^2R$ heating in the coils and eddy currents in the core. Fixes: closed iron
        core, thick copper wire, laminated core.
  \item Line loss $P_\mathrm{loss} = I^2R$ (not on the sheet). For a fixed power, a higher
        voltage gives a smaller current, so much less loss.
  \item Grid: step up at the power station; step down in stages to \qty{230}{V} for safety.
\end{itemize}
\end{summarybox}
